%%%%%%%%%%%%%%%%
%%% num 27 %%%
%%%%%%%%%%%%%%%%

\Chapter{27}

\Verse{1}
{
	\hP{וַתִּקְרַבְנָה בְּנוֹת צְלפְחָד בֶּן חֵפֶר בֶּן גִּלְעָד בֶּן מָכִיר בֶּן מְנַשֶּׁה לְמִשְׁפְּחֹת מְנַשֶּׁה בֶן יוֹסֵף וְאֵלֶּה שְׁמוֹת בְּנֹתָיו מַחְלָה נֹעָה וְחגְלָה וּמִלְכָּה וְתִרְצָה׃}
}
{
	\eP{
		27 And the daughters of Zelophehad son of Hepher son of\\
		Gilead son of Machir son of Manasseh, of the families of\\
		Manasseh son of Joseph, came forward. And these are his\\
		daughters’ names: Mahlah, Noah, and Hoglah and Milcah and\\
		Tirzah. {[image]}\\
	}
}
{
}

\Verse{2}
{
	\hP{וַתַּעֲמֹדְנָה לִפְנֵי מֹשֶׁה וְלִפְנֵי אֶלְעָזָר הַכֹּהֵן וְלִפְנֵי הַנְּשִׂיאִם וְכל הָעֵדָה פֶּתַח אֹהֶל מוֹעֵד לֵאמֹר׃}
}
{
	\eP{
		And they stood in front of Moses and in front of Eleazar,\\
		the priest, and in front of the chieftains and all the\\
		congregation at the entrance of the Tent of Meeting,\\
		saying, {[image]}\\
	}
}
{
}

\Verse{3}
{
	\hP{אָבִינוּ מֵת בַּמִּדְבָּר וְהוּא לֹא הָיָה בְּתוֹךְ הָעֵדָה הַנּוֹעָדִים עַל יְהֹוָה בַּעֲדַת קֹרַח כִּי בְחֶטְאוֹ מֵת וּבָנִים לֹא הָיוּ לוֹ׃}
}
{
	\eP{
		“Our father died in the wilderness. And he wasn’t among\\
		the congregation who were gathered against YHWH: in\\
		Korah’s congregation. Rather, he died through his own sin.\\
		And he had no sons. {[image]}\\
	}
}
{
}

\Verse{4}
{
	\hP{לָמָּה יִגָּרַע שֵׁם אָבִינוּ מִתּוֹךְ מִשְׁפַּחְתּוֹ כִּי אֵין לוֹ בֵּן תְּנָה לָּנוּ אֲחֻזָּה בְּתוֹךְ אֲחֵי אָבִינוּ׃}
}
{
	\eP{
		Why should our father’s name be subtracted from among his\\
		family because he didn’t have a son? Give us a possession\\
		among our father’s brothers.” {[image]}\\
	}
}
{
%	\footnote{
%		Why should our father’s name be subtracted. They do not make their case
%		in terms of their own rights or of women’s rights in general. They put
%		it in terms of their father, that his name would be eliminated from the
%		territorial holdings of Israel. The issue is the one that arises in the
%		law of the jubilee. Land is unalienable. It is to be distributed among
%		the Israelites when they arrive in the land. If a man ever loses his
%		land, it is returned to him in the jubilee year. If he has died by then,
%		it is returned to his sons. But the property in his name stays in the
%		family forever. Zelophehad’s daughters now protest that the property
%		that would have gone to their father will go to his brothers instead,
%		and the family holding in his name will be lost.
%	}
}

\Verse{5}
{
	\hP{וַיַּקְרֵב מֹשֶׁה אֶת מִשְׁפָּטָן לִפְנֵי יְהֹוָה׃}
}
{
	\eP{
		And Moses brought their case forward in front of YHWH.\\
		{[image]}\\
	}
}
{
}

\Verse{6}
{
	\hP{וַיֹּאמֶר יְהֹוָה אֶל מֹשֶׁה לֵּאמֹר׃}
}
{
	\eP{And YHWH said to Moses, saying, [image]}
}
{
}

\Verse{7}
{
	\hP{כֵּן בְּנוֹת צְלפְחָד דֹּבְרֹת נָתֹן תִּתֵּן לָהֶם אֲחֻזַּת נַחֲלָה בְּתוֹךְ אֲחֵי אֲבִיהֶם וְהַעֲבַרְתָּ אֶת נַחֲלַת אֲבִיהֶן לָהֶן׃}
}
{
	\eP{
		“Zelophehad’s daughters speak right. You shall give them a\\
		possession for a legacy among their father’s brothers, and\\
		you shall pass their father’s legacy to them. {[image]}\\
	}
}
{
%	\footnote{
%		Zelophehad’s daughters speak right. See the comment on Num 36:5.
%	}
}

\Verse{8}
{
	\hP{וְאֶל בְּנֵי יִשְׂרָאֵל תְּדַבֵּר לֵאמֹר אִישׁ כִּי יָמוּת וּבֵן אֵין לוֹ וְהַעֲבַרְתֶּם אֶת נַחֲלָתוֹ לְבִתּוֹ׃}
}
{
	\eP{
		And you shall speak to the children of Israel, saying, ‘A\\
		man who will die and not have a son: you shall pass his\\
		legacy to his daughter. {[image]}\\
	}
}
{
%	\footnote{
%		A man who will die and not have a son: you shall pass his legacy to his
%		daughter. This judgment is an important step in the development of
%		women’s rights, but its message is mixed. On one hand, it says that
%		women can inherit property, and their right of inheritance precedes the
%		rights of their father’s male siblings or any other male relatives who
%		are more distantly related to their father than the women themselves. On
%		the other hand, this applies only if their father had no sons. If a
%		father has even one son and ten daughters, the son inherits the family
%		land. The daughters are dependent on that brother or on their husbands
%		for property. (This matter of women’s inheritance will become even more
%		complex in the last chapter of Numbers; see the comments on Num 36:3;
%		36:5; and 36:6.) There is little point in debating whether this step
%		means that the Torah is supportive of women, on the grounds that it
%		provides for them to inherit, or whether it means that the Torah is
%		unfair to women, on the grounds that sons still precede daughters. The
%		fact is that social transformations take time: generations, centuries,
%		even millennia. The Torah does not command a revolution in the status of
%		women. It provides for steps such as this one, which in the short run
%		established that women do have rights, and which ultimately participated
%		in the development of women’s rights more generally. We can praise the
%		Bible for how far it went, or we can be critical that it did not go
%		farther. But we would do better to examine how far it went in its age,
%		and how much this contributed to the transformation in the balance
%		between men and women in the millennia that followed. The larger point
%		is the same that I made with regard to slavery: The Torah does not
%		forbid it and attempt to bring it to an end overnight. It rather gives
%		laws of treatment of slaves—which involved granting respect, rights, and
%		compassion for slaves. And this eventually undermined slavery as an
%		institution. The diminution of slavery and the increase of women’s
%		rights are two of the major developments of the past century. The
%		Torah’s laws played an early and determinative part in birthing and
%		nurturing both of these revolutions.
%	}
}

\Verse{9}
{
	\hP{וְאִם אֵין לוֹ בַּת וּנְתַתֶּם אֶת נַחֲלָתוֹ לְאֶחָיו׃}
}
{
	\eP{
		And if he does not have a daughter, then you shall give\\
		his legacy to his brothers. {[image]}\\
	}
}
{
}

\Verse{10}
{
	\hP{וְאִם אֵין לוֹ אַחִים וּנְתַתֶּם אֶת נַחֲלָתוֹ לַאֲחֵי אָבִיו׃}
}
{
	\eP{
		And if he does not have brothers, then you shall give his\\
		legacy to his father’s brothers. {[image]}\\
	}
}
{
}

\Verse{11}
{
	\hP{וְאִם אֵין אַחִים לְאָבִיו וּנְתַתֶּם אֶת נַחֲלָתוֹ לִשְׁאֵרוֹ הַקָּרֹב אֵלָיו מִמִּשְׁפַּחְתּוֹ וְיָרַשׁ אֹתָהּ וְהָיְתָה לִבְנֵי יִשְׂרָאֵל לְחֻקַּת מִשְׁפָּט כַּאֲשֶׁר צִוָּה יְהֹוָה אֶת מֹשֶׁה׃}
}
{
	\eP{
		And if there are no brothers of his father, then you shall\\
		give his legacy to his relative who is closest to him of\\
		his family, and he shall possess it.’ And it shall become\\
		a law of judgment for the children of Israel, as YHWH\\
		commanded Moses.” {[image]}\\
	}
}
{
%	\footnote{
%		law of judgment. This term is used twice in the Torah: here for cases of
%		inheritance, and later for cases involving cities of refuge (Num 35:29).
%		In both occurrences it involves the stating of a legal procedure for
%		deciding cases—civil procedure in this instance, and criminal procedure
%		in the later instance—which judges will then apply to the particular
%		cases before them.
%	}
}

\Verse{12}
{
	\hP{וַיֹּאמֶר יְהֹוָה אֶל מֹשֶׁה עֲלֵה אֶל הַר הָעֲבָרִים הַזֶּה וּרְאֵה אֶת הָאָרֶץ אֲשֶׁר נָתַתִּי לִבְנֵי יִשְׂרָאֵל׃}
}
{
	\eP{
		And YHWH said to Moses, “Go up to this mountain of Abarim\\
		and see the land that I’ve given to the children of\\
		Israel. {[image]}\\
	}
}
{
}

\Verse{13}
{
	\hP{וְרָאִיתָה אֹתָהּ וְנֶאֱסַפְתָּ אֶל עַמֶּיךָ גַּם אָתָּה כַּאֲשֶׁר נֶאֱסַף אַהֲרֹן אָחִיךָ׃}
}
{
	\eP{
		And you’ll see it, and then you’ll be gathered to your\\
		people, you as well, as Aaron, your brother, was gathered,\\
		{[image]}\\
	}
}
{
}

\Verse{14}
{
	\hP{כַּאֲשֶׁר מְרִיתֶם פִּי בְּמִדְבַּר צִן בִּמְרִיבַת הָעֵדָה לְהַקְדִּישֵׁנִי בַמַּיִם לְעֵינֵיהֶם הֵם מֵי מְרִיבַת קָדֵשׁ מִדְבַּר צִן׃}
}
{
	\eP{
		because you rebelled against my word in the wilderness of\\
		Zin in the congregation’s quarrel to make me holy with\\
		water before their eyes. That is the water of Meribah of\\
		Kadesh at the wilderness of Zin.” {[image]}\\
	}
}
{
}

\Verse{15}
{
	\hP{וַיְדַבֵּר מֹשֶׁה אֶל יְהֹוָה לֵאמֹר׃}
}
{
	\eP{And Moses spoke to YHWH, saying, [image]}
}
{
%	\footnote{
%		Moses spoke to YHWH. This is the last time that Moses is quoted as
%		saying anything to God. In Deuteronomy Moses will report to the people
%		his past conversations with God, but these are the last words of Moses
%		to God in the Torah. And his last concern is: the people—that they
%		should have a leader. And note that all of this draws us back to Moses’
%		first meeting with God, at the burning bush. God begins the charge to
%		Moses there by expressing concern for the people: “I’ve seen the
%		degradation of my people” (Exod 3:7). So God’s first charge to Moses and
%		Moses’ last request of God are both about the good of the people. And
%		the connection to Moses’ origins is underscored by Moses’ choice of
%		metaphor here. The episode at the bush begins: “Moses had been
%		shepherding the flock” (Exod 3:1); and now Moses asks for a leader so
%		the people won’t be “like sheep that don’t have a shepherd” (Num 27:17).
%	}
}

\Verse{16}
{
	\hP{יִפְקֹד יְהֹוָה אֱלֹהֵי הָרוּחֹת לְכל בָּשָׂר אִישׁ עַל הָעֵדָה׃}
}
{
	\eP{
		“Let YHWH, God of the spirits of all flesh, appoint a man\\
		over the congregation {[image]}\\
	}
}
{
}

\Verse{17}
{
	\hP{אֲשֶׁר יֵצֵא לִפְנֵיהֶם וַאֲשֶׁר יָבֹא לִפְנֵיהֶם וַאֲשֶׁר יוֹצִיאֵם וַאֲשֶׁר יְבִיאֵם וְלֹא תִהְיֶה עֲדַת יְהֹוָה כַּצֹּאן אֲשֶׁר אֵין לָהֶם רֹעֶה׃}
}
{
	\eP{
		who will go out in front of them and who will come in in\\
		front of them and who will bring them out and who will\\
		bring them in, so YHWH’s congregation won’t be like sheep\\
		that don’t have a shepherd.” {[image]}\\
	}
}
{
}

\Verse{18}
{
	\hP{וַיֹּאמֶר יְהֹוָה אֶל מֹשֶׁה קַח לְךָ אֶת יְהוֹשֻׁעַ בִּן נוּן אִישׁ אֲשֶׁר רוּחַ בּוֹ וְסָמַכְתָּ אֶת יָדְךָ עָלָיו׃}
}
{
	\eP{
		And YHWH said to Moses, “Take Joshua son of Nun, a man\\
		with spirit in him, and lay your hand on him. {[image]}\\
	}
}
{
}

\Verse{19}
{
	\hP{וְהַעֲמַדְתָּ אֹתוֹ לִפְנֵי אֶלְעָזָר הַכֹּהֵן וְלִפְנֵי כּל הָעֵדָה וְצִוִּיתָה אֹתוֹ לְעֵינֵיהֶם׃}
}
{
	\eP{
		And you shall stand him in front of Eleazar, the priest,\\
		and in front of all the congregation, and you shall\\
		command him before their eyes. {[image]}\\
	}
}
{
}

\Verse{20}
{
	\hP{וְנָתַתָּה מֵהוֹדְךָ עָלָיו לְמַעַן יִשְׁמְעוּ כּל עֲדַת בְּנֵי יִשְׂרָאֵל׃}
}
{
	\eP{
		And you shall put some of your eminence on him so that all\\
		the congregation of the children of Israel will hear.\\
		{[image]}\\
	}
}
{
}

\Verse{21}
{
	\hP{וְלִפְנֵי אֶלְעָזָר הַכֹּהֵן יַעֲמֹד וְשָׁאַל לוֹ בְּמִשְׁפַּט הָאוּרִים לִפְנֵי יְהֹוָה עַל פִּיו יֵצְאוּ וְעַל פִּיו יָבֹאוּ הוּא וְכל בְּנֵי יִשְׂרָאֵל אִתּוֹ וְכל הָעֵדָה׃}
}
{
	\eP{
		And he shall stand in front of Eleazar, the priest; and he\\
		shall ask him in judgment of the Urim in front of YHWH. By\\
		his word they shall go out, and by his word they shall\\
		come in, he and all the children of Israel with him and\\
		all the congregation.” {[image]}\\
	}
}
{
}

\Verse{22}
{
	\hP{וַיַּעַשׂ מֹשֶׁה כַּאֲשֶׁר צִוָּה יְהֹוָה אֹתוֹ וַיִּקַּח אֶת יְהוֹשֻׁעַ וַיַּעֲמִדֵהוּ לִפְנֵי אֶלְעָזָר הַכֹּהֵן וְלִפְנֵי כּל הָעֵדָה׃}
}
{
	\eP{
		And Moses did as YHWH commanded him, and he took Joshua\\
		and stood him in front of Eleazar, the priest, and in\\
		front of all the congregation. {[image]}\\
	}
}
{
}

\Verse{23}
{
	\hP{וַיִּסְמֹךְ אֶת יָדָיו עָלָיו וַיְצַוֵּהוּ כַּאֲשֶׁר דִּבֶּר יְהֹוָה בְּיַד מֹשֶׁה׃}
}
{
	\eP{
		And he laid his hands on him and commanded him, as YHWH\\
		had spoken by Moses’ hand. {[image]}\\
	}
}
{
%	\footnote{
%		he laid his hands. God used the singular: “Lay your hand.” But Moses
%		lays both hands on Joshua. It is as if he wants to make the appointment
%		of his successor as visible and unmistakable as possible. And he
%		indicates in this way that this appointment is his own choice as well.
%		He fulfills the divine command with one hand. The second hand is his own
%		affirmation of Joshua, which bolsters Joshua’s appointment and preempts
%		anyone from criticizing Joshua later as not having been Moses’ choice.
%		The successor to a great leader is always in a vulnerable position, and
%		it is a gracious act by the great leader to support that successor.
%	}
}
